% Compile with pdfLaTeX.
% Keep PGF-02-R39_G10_Global_Economics_Term_1_2026-2027_logo.jpeg
% in the same directory as this file to display the original school logo.
\documentclass[10pt,letterpaper]{article}

\usepackage[landscape,margin=0.5in,bottom=0.62in]{geometry}
\usepackage[T1]{fontenc}
\usepackage[utf8]{inputenc}
\usepackage[scaled=0.94]{helvet}
\renewcommand{\familydefault}{\sfdefault}
\usepackage{microtype}
\usepackage{graphicx}
\usepackage[table]{xcolor}
\usepackage{array,tabularx,multirow}
\usepackage[most]{tcolorbox}
\usepackage{ragged2e}
\usepackage{fancyhdr}

\definecolor{sectionblue}{HTML}{4BACC6}
\definecolor{borderblue}{HTML}{4F81BD}
\definecolor{highlightyellow}{HTML}{FFFF00}

\setlength{\parindent}{0pt}
\setlength{\parskip}{0pt}
\setlength{\tabcolsep}{8pt}
\renewcommand{\arraystretch}{1.12}

\pagestyle{fancy}
\fancyhf{}
\renewcommand{\headrulewidth}{0pt}
\renewcommand{\footrulewidth}{0pt}
\fancyfoot[L]{\colorbox{highlightyellow}{\strut\rmfamily\fontsize{10}{11}\selectfont Versión 08}}

\newcommand{\schoollogo}{%
  \IfFileExists{PGF-02-R39_G10_Global_Economics_Term_1_2026-2027_logo.jpeg}{%
    \includegraphics[width=0.72in]{PGF-02-R39_G10_Global_Economics_Term_1_2026-2027_logo.jpeg}%
  }{%
    \fbox{\parbox[c][0.68in][c]{0.68in}{\centering\scriptsize SAN\\BARTOLOMÉ\\[1pt]JHS\\[1pt]LA MERCED}}%
  }%
}

\newcommand{\sectionbar}[1]{%
  \begin{tcolorbox}[
    enhanced,
    colback=sectionblue,
    colframe=sectionblue,
    boxrule=0pt,
    sharp corners,
    left=0.32in,
    right=0.08in,
    top=3.5pt,
    bottom=3.5pt,
    before skip=0pt,
    after skip=0pt,
    borderline north={1.15pt}{0pt}{black},
    borderline south={1.15pt}{0pt}{black}
  ]
  \color{white}\bfseries\fontsize{12}{14}\selectfont #1
  \end{tcolorbox}%
}

\newtcolorbox{bluebox}[1][]{
  enhanced,
  colback=white,
  colframe=borderblue,
  boxrule=1.15pt,
  sharp corners,
  left=6pt,
  right=6pt,
  top=5pt,
  bottom=5pt,
  before skip=0pt,
  after skip=0pt,
  #1
}

\newtcolorbox{plainbox}[1][]{
  enhanced,
  colback=white,
  colframe=black,
  boxrule=0.45pt,
  sharp corners,
  left=6pt,
  right=6pt,
  top=5pt,
  bottom=5pt,
  before skip=0pt,
  after skip=0pt,
  #1
}

\newcommand{\criterion}[2]{\textbf{C#1.} #2\par}
\newcommand{\field}[2]{\textbf{#1:} #2\par}
\newcommand{\yellowmark}[1]{\colorbox{highlightyellow}{\strut\textcolor{white}{#1}}}

\begin{document}
\fontsize{9.5}{12}\selectfont

% -----------------------------------------------------------------------------
% Institutional header
% -----------------------------------------------------------------------------
\noindent
\begin{tabularx}{\textwidth}{|>{\centering\arraybackslash}m{0.95in}|>{\centering\arraybackslash}X|>{\centering\arraybackslash}m{1.18in}|}
\hline
\begin{minipage}[c][0.86in][c]{0.82in}\centering\schoollogo\end{minipage} &
  {\bfseries\fontsize{12}{14}\selectfont TEACHING-LEARNING PLAN} &
  \begin{minipage}[c][0.86in][c]{1.08in}
    \centering
    \colorbox{highlightyellow}{\bfseries\fontsize{10.5}{12}\selectfont PGF-02-R39}\par
    \vspace{5pt}\hrule\vspace{4pt}
    \textbf{DATE}\par
    \vspace{2pt}
    August 12 --\par November 5, 2026
  \end{minipage}
\\
\hline
\end{tabularx}

\vspace{0.23in}

\sectionbar{1.\quad Identification}

\vspace{0.23in}

\begin{plainbox}
  \field{Area of knowledge}{Global Economics}
  \field{Teacher's name (or Teachers' names)}{Nicolas Lopez}
  \field{Grade}{10}
  \field{Term}{1 (2026-2027)}
\end{plainbox}

\vspace{0.65in}

\sectionbar{2.\quad Learning Objective}

\vspace{0.22in}

\begin{minipage}[t]{0.28\textwidth}
  \centering\bfseries\fontsize{12}{14}\selectfont LEARNING OBJECTIVE
\end{minipage}%
\hfill
\begin{minipage}[t]{0.70\textwidth}
  \centering\bfseries\fontsize{12}{14}\selectfont ASSESSMENT CRITERIA
\end{minipage}

\vspace{3pt}

\begin{minipage}[t]{0.28\textwidth}
  \begin{bluebox}[height=2.82in,valign=top]
    \fontsize{9.5}{12.4}\selectfont
    Analyze economic decisions under uncertainty by formulating alternatives,
    events, states and consequences; constructing payoff tables from direct
    payoffs or revenue and cost information; applying and comparing decision
    criteria with and without probabilities; and communicating a justified
    recommendation through tables, decision trees and project evidence.
  \end{bluebox}
\end{minipage}%
\hfill
\begin{minipage}[t]{0.70\textwidth}
  \begin{bluebox}[height=2.82in,valign=top]
    \fontsize{8.05}{10.7}\selectfont
    \criterion{1}{Characterize a decision under uncertainty by identifying its context, alternatives, event, states of nature and consequences.}
    \criterion{2}{Interpret and organize alternatives, events, states, consequences and supplied payoffs in a payoff table.}
    \criterion{3}{Calculate revenue, variable and fixed costs, quantities and profit, and construct the corresponding payoff table.}
    \criterion{4}{Apply Maximax and Maximin to select and justify optimistic and conservative alternatives without probabilities.}
    \criterion{5}{Construct opportunity-loss tables and apply Minimax Regret.}
    \criterion{6}{Use probabilities to calculate expected monetary value and select the preferred alternative.}
    \criterion{7}{Analyze the sensitivity of expected-value recommendations to changes in probabilities.}
    \criterion{8}{Compare Bayes/expected value, Maximin and Minimax Regret on common decision problems.}
    \criterion{9}{Represent staged decisions and uncertainty with decision trees.}
    \criterion{10}{Conduct and communicate an end-to-end decision analysis that integrates models, criteria, results and a justified recommendation.}
  \end{bluebox}
\end{minipage}

\newpage

% -----------------------------------------------------------------------------
% Conceptual standards, context, and experience
% -----------------------------------------------------------------------------
\sectionbar{3. Conceptual standards of the disciplinary knowledge that supports the narrative's development.}

\vspace{0.23in}

\begin{bluebox}
  \begin{center}
  \begin{minipage}{0.89\textwidth}
    \textbf{Core disciplinary knowledge:} decisions under uncertainty;
    decision-makers; alternatives; events; states of nature; consequences;
    direct and calculated payoffs; revenue, quantities, variable and fixed
    costs, profit and payoff tables.
  \end{minipage}
  \end{center}
  \textbf{Analytical standards:} Maximax, Maximin, opportunity loss and Minimax
  Regret; probabilities, Bayes/expected monetary value and sensitivity analysis;
  decision trees; integrated comparison, interpretation and justification of
  recommendations.
\end{bluebox}

\vspace{0.24in}

\sectionbar{4.\quad Pre-lesson and Context:
  \yellowmark{\fontsize{10.7}{12}\selectfont Pedagogical actions planned by the teacher to motivate students for the term work proposal}\par
  \hspace*{0.49in}\yellowmark{\fontsize{10.7}{12}\selectfont (presentation of scenarios, narratives, learning objectives and assessment criteria)}}

\vspace{0.22in}

\begin{bluebox}
  Begin with familiar uncertain situations from retail, transport, energy and
  school operations. Use the Term 1 survey and a short initial scenario to
  elicit students' decision processes, then introduce the term learning
  objective, C1-C10 criteria, the three learning-evidence blocks and the
  payoff-table structure. Share materials and instructions through OPEN/Teams.
\end{bluebox}

\vspace{0.23in}

\sectionbar{5.\quad Experience:
  \yellowmark{\fontsize{10.7}{12}\selectfont pedagogical actions that establish the learning objective according to the assessment criteria.}\par
  \hspace*{0.49in}\yellowmark{\fontsize{10.7}{12}\selectfont Links to virtual resources or ways used to share experiences (OPEN, TEAMS, etc.) should be included.}}

\vspace{0.22in}

\renewcommand{\arraystretch}{1.10}
\noindent
\begin{tabularx}{\textwidth}{|>{\RaggedRight\arraybackslash}p{2.42in}|>{\RaggedRight\arraybackslash}X|}
\hline
\textbf{PEDAGOGICAL ACTION: N1}\par
\textbf{C1-C4}\par
\textit{First Learning Evidence}
&
\textbf{Learning sequence:} Develop Practice Activities 1-3 through teacher
modeling, guided analysis and individual or paired work. Students identify
decision elements, organize and calculate payoffs, construct payoff tables, and
apply Maximax and Maximin.\par
\textbf{Learning evidence:} Written exam based on the practice activities,
assessing C1-C4. Use the criterion-specific catch-up materials for feedback and
remediation.
\\
\hline
\textbf{PEDAGOGICAL ACTION: N2}\par
\textbf{C5-C7}\par
\textit{Second Learning Evidence}
&
\textbf{Learning sequence:} Complete the regret component of Practice Activity
3 and Practice Activity 4. Students construct opportunity-loss tables, apply
Minimax Regret, calculate expected monetary value, and test how recommendations
change as probabilities vary.\par
\textbf{Learning evidence:} Written exam based on the practice activities,
assessing C5-C7. Follow with targeted correction and criterion-specific catch-up
work when needed.
\\
\hline
\end{tabularx}

\newpage

\noindent
\begin{tabularx}{\textwidth}{|>{\RaggedRight\arraybackslash}p{2.42in}|>{\RaggedRight\arraybackslash}X|}
\hline
\textbf{PEDAGOGICAL ACTION: N3}\par
\textbf{C8-C10}\par
\textit{Third Learning Evidence}
&
\textbf{Learning sequence:} Use the C8-C10 extension practice to compare
Bayes/expected value, Maximin and Minimax Regret on common data; construct
decision trees; and complete an end-to-end analysis of an economic decision
under uncertainty.\par
\textbf{Learning evidence:} Written exam based on the practice activities,
complemented by project work in which students model a case, apply the relevant
methods, interpret results and justify a final recommendation.
\\
\hline
\end{tabularx}

\vspace{0.23in}

{\bfseries\fontsize{12}{14}\selectfont References:}

\vspace{3pt}

\begin{bluebox}[height=1.95in,valign=top]
  \textbf{Course map:} README.md - Grade 10 Global Economics, Term 1.\par
  \textbf{Practice Activity 1:} \textnormal{\detokenize{G10_T1_L1_C1C2_practice_activity1.tex}} - decision elements and payoff tables.\par
  \textbf{Practice Activity 2:} \textnormal{\detokenize{G10_T1_L1_C1C3_practice_activity2.tex}} - revenue, costs, profit and payoff tables.\par
  \textbf{Practice Activity 3:} \textnormal{\detokenize{G10_T1_L2_C4C5_practice_activity3.tex}} - Maximax, Maximin and Minimax Regret.\par
  \textbf{Practice Activity 4:} \textnormal{\detokenize{G10_T1_L2_C6C7_practice_activity4.tex}} - expected value and sensitivity analysis.\par
  \textbf{Extension material:} \textnormal{\detokenize{G10_T1_L3_MATERIAL.tex}} - method comparison, decision trees and integrated C8-C10 analysis.
\end{bluebox}

\end{document}
